\documentclass[a4paper,twoside]{article}
\usepackage{blindtext}  
\usepackage{geometry}

% Chinese support
\usepackage[UTF8, scheme = plain]{ctex}

% Page margin layout
\geometry{left=2.3cm,right=2cm,top=2.5cm,bottom=2.0cm}


\usepackage{listings}
\usepackage{xcolor}
\usepackage{geometry}
\usepackage{amsmath}
\usepackage{float}
\usepackage{hyperref}

\usepackage{graphics}
\usepackage{graphicx}
\usepackage{subfigure}
\usepackage{epsfig}
\usepackage{float}

\usepackage{algorithm}
\usepackage[noend]{algpseudocode}

\usepackage{booktabs}
\usepackage{threeparttable}
\usepackage{longtable}
\usepackage{listings}
\usepackage{tikz}
\usepackage{multicol}

% cite package, to clean up citations in the main text. Do not remove.
\usepackage{cite}

\usepackage{color,xcolor}

%% The amssymb package provides various useful mathematical symbols
\usepackage{amssymb}
%% The amsthm package provides extended theorem environments
\usepackage{amsthm}
\usepackage{amsfonts}
\usepackage{enumerate}
\usepackage{enumitem}
\usepackage{listings}

\usepackage{indentfirst}
\setlength{\parindent}{2em} % Make two letter space in the first paragraph
\usepackage{setspace}
\linespread{1.5} % Line spacing setting
\usepackage{siunitx}
\setlength{\parskip}{0.5em} % Paragraph spacing setting

% \usepackage[contents =22920202204622, scale = 10, color = black, angle = 50, opacity = .10]{background}

\renewcommand{\figurename}{图}
\renewcommand{\lstlistingname}{代码} 
\renewcommand{\tablename}{表格}
\renewcommand{\contentsname}{目录}
\floatname{algorithm}{算法}

\graphicspath{ {images/} }

%%%%%%%%%%%%%
\newcommand{\StudentNumber}{22920202204622}  % Fill your student number here
\newcommand{\StudentName}{熊恪峥}  % Replace your name here
\newcommand{\PaperTitle}{作业（十）}  % Change your paper title here
\newcommand{\PaperType}{算法设计与分析} % Replace the type of your report here
\newcommand{\Date}{2022年5月25日}
\newcommand{\College}{信息学院}
\newcommand{\CourseName}{算法设计与分析}
%%%%%%%%%%%%%

%% Page header and footer setting
\usepackage{fancyhdr}
\usepackage{lastpage}
\pagestyle{fancy}
\fancyhf{}
% This requires the document to be twoside
\fancyhead[LO]{\texttt{\StudentName }}
\fancyhead[LE]{\texttt{\StudentNumber}}
\fancyhead[C]{\texttt{\PaperTitle }}
\fancyhead[R]{\texttt{第{\thepage}页，共\pageref*{LastPage}页}}


\title{\PaperTitle}
\author{\StudentName}
\date{\Date}

\lstset{
	basicstyle          =   \sffamily,          % 基本代码风格
	keywordstyle        =   \bfseries,          % 关键字风格
	commentstyle        =   \rmfamily\itshape,  % 注释的风格，斜体
	stringstyle         =   \ttfamily,  % 字符串风格
	flexiblecolumns,                % 别问为什么，加上这个
	numbers             =   left,   % 行号的位置在左边
	showspaces          =   false,  % 是否显示空格，显示了有点乱，所以不现实了
	numberstyle         =   \zihao{-5}\ttfamily,    % 行号的样式，小五号，tt等宽字体
	showstringspaces    =   false,
	captionpos          =   t,      % 这段代码的名字所呈现的位置，t指的是top上面
	frame               =   lrtb,   % 显示边框
}

\lstdefinestyle{PythonStyle}{
	language        =   Python, % 语言选Python
	basicstyle      =   \zihao{-5}\ttfamily,
	numberstyle     =   \zihao{-5}\ttfamily,
	keywordstyle    =   \color{blue},
	keywordstyle    =   [2] \color{teal},
	stringstyle     =   \color{magenta},
	commentstyle    =   \color{red}\ttfamily,
	breaklines      =   true,   % 自动换行，建议不要写太长的行
	columns         =   fixed,  % 如果不加这一句，字间距就不固定，很丑，必须加
	basewidth       =   0.5em,
}

\lstdefinestyle{CppStyle}{
	language        =   c++,
	basicstyle      =   \zihao{-5}\ttfamily,
	numberstyle     =   \zihao{-5}\ttfamily,
	keywordstyle    =   \color{blue},
	keywordstyle    =   [2] \color{teal},
	stringstyle     =   \color{magenta},
	commentstyle    =   \color{red}\ttfamily,
	breaklines      =   true,   % 自动换行，建议不要写太长的行
	columns         =   fixed,  % 如果不加这一句，字间距就不固定，很丑，必须加
	basewidth       =   0.5em,
}

\algnewcommand\algorithmicinput{\textbf{Input:}}
\algnewcommand\algorithmicoutput{\textbf{Output:}}
\algnewcommand\Input{\item[\algorithmicinput]}%
\algnewcommand\Output{\item[\algorithmicoutput]}%

\usetikzlibrary{positioning, shapes.geometric}

\begin{document}
	
%%%%%%%%%%%%%%%%%%%%%%%%%%%%%%%%%%%%%%%%%%%%
\makeatletter % change default title style
\renewcommand*\maketitle{%
	\begin{center} 
		\bfseries  % title 
		{\LARGE \@title \par}  % LARGE typesetting
		\vskip 1em  %  margin 1em
		{\global\let\author\@empty}  % no author information
		{\global\let\date\@empty}  % no date
		\thispagestyle{empty}   %  empty page style
	\end{center}%
	\setcounter{footnote}{0}%
}
\makeatother
%%%%%%%%%%%%%%%%%%%%%%%%%%%%%%%%%%%%%%%%%%%%
	
	
\thispagestyle{empty}

\vspace*{1cm}

\begin{figure}[h]
	\centering
	\includegraphics[width=4.0cm]{logo.png}
\end{figure}

\vspace*{1cm}

\begin{center}
	\Huge{\textbf{\PaperType}}
	
	\Large{\PaperTitle}
\end{center}

\vspace*{1cm}

\begin{table}[h]
	\centering	
	\begin{Large}
		\renewcommand{\arraystretch}{1.5}
		\begin{tabular}{p{3cm} p{5cm}<{\centering}}
			姓\qquad 名 & \StudentName  \\
			\hline
			学\qquad号 & \StudentNumber \\
			\hline
			日\qquad期 & \Date  \\
			\hline
			学\qquad院 & \College  \\
			\hline
			课程名称 & \CourseName  \\
			\hline
		\end{tabular}
	\end{Large}
\end{table}

\newpage

\title{
	\Large{\textcolor{black}{\PaperTitle}}
}
	
	
\maketitle
	
\tableofcontents
 
\newpage
\setcounter{page}{1}

\begin{spacing}{1.2}

\section{题11.2}

\begin{proof}
	先证充分性。

	若$LongestPathLength$为多项式时间可解的问题，设时间复杂度为$\mathcal{O}(P(n))$，则
	解决$LongestPath$问题可等价于判定
	\[
		LongestPathLength(G,u,v) \ne \infty
	\]
	这个方法的时间复杂度就是$\mathcal{O}(P(n)+1)=\mathcal{O}(P(n))$，因此$LongestPath$问题多项式
	时间可解。

	再证必要性。

	若$LongestPath$为多项式时间可解的问题，设时间复杂度为$\mathcal{O}(P(n))$，则算法\ref{algo:lp}
	可以解决$LongestPathLength$问题。
	\begin{algorithm}[htbp]
		\caption{解决$LongestPathLength$}
		\label{algo:lp}
		\begin{algorithmic}[1]
			\Procedure{LongestPathLength}{$p,k,n$}
				\For{$i=1\to |\mathcal{V}|$}
					\If{$\mathop{not}$ \Call{LongestPath}{G,u,v,i}}
						\State \Return $i-1$
					\EndIf
				\EndFor
			\EndProcedure
		\end{algorithmic}
	\end{algorithm}	

	$LongestPathLength$的内层循环次数不超过$|\mathcal{V}|$，它的时间复杂度是$\mathcal{O}(|\mathcal{V}|P(n))$，
	因此是多项式时间可解的。

\end{proof}

\section{题11.3}

假设TSP问题是多项式时间可解的，求出回路的权值只需要在算法给出的顶点序列中累加对应边的权值。
如算法\ref{algo:tsplen}。

\begin{algorithm}[htbp]
	\caption{解决$TSPLength$}
	\label{algo:tsplen}
	\begin{algorithmic}[1]
		\Procedure{TSPLength}{$G$}
			\State $\mathcal{V}=TSP(G)$
			\State $len\gets 0$
			\For{$i=1\to |\mathcal{V}|$}
				\State $len \gets len+w[\mathcal{V}[i],\mathcal{V}[i-1]]$
			\EndFor
			\State $len\gets len+w[\mathcal{V}[|\mathcal{V}|-1],0]$
			\State \Return $len$
		\EndProcedure
	\end{algorithmic}
\end{algorithm}	

\section{题11.4}

\textbf{优化问题：}找出一条从$u$到$u$的简单回路，使得这条回路的权值和$w$最大。

\textbf{判定问题：}是否存在一条从$u$到$u$的简单回路，这条回路的权值和至少为$w$。

它的语言如\eqref{eqn:langlongest}。

\begin{equation}
	\label{eqn:langlongest}
	\begin{aligned}
		LongestPath=\{ <G,u,w>: G(V,E)\mbox{是无向图}, u\in V, \mbox{存在从}u\mbox{开始的回路，它的权值和不少于}w \}
	\end{aligned}
\end{equation}

\section{题11.7}

若$HamCycle$多项式时间可解，设时间复杂度是$\mathcal{O}(P(n))$，注意到哈密顿回路的每个顶点都连着两条边。
可以对顶点$v$枚举所有与之相连的边集合$E_v$，找出两条边使得$G'(V,(E-E_v)\cup \{e_i,e_j\})$是哈密顿图。并
沿着这两条边的另一顶点枚举，直到对所有顶点枚举完成。所有找出的边就是哈密顿回路的边。
这个过程的时间复杂度是$\mathcal{O}(|V|^2P(n))$，也是多项式的。之后就可以沿着这些边列出所有的顶点。


\section{题11.16}

先证01整数规划问题是NP的。
\begin{proof}
	01整数规划问题的判定问题是\emph{给定矩阵$A_{mn}$、$x$和$b$，是否有$Ax\le b$}。
	
	给定矩阵$A_{mn}$、$x$和$b$，判断是否$Ax\le b$需要进行一次矩阵乘法和$n$次整数比较。
	可以在多项式时间完成。

	因此01整数规划问题是NP的。
\end{proof}

然后证明01整数规划问题是NP-Hard的。要证01整数规划问题是NP-Hard的，就要证明存在一个NP-Complete问题，它可以多项式时间
约简到01整数规划问题。可以证明01INT$\le_p$3-CNF-SAT
\begin{proof}
由于$x_i\in \{0,1\}$,有
\[
	0\le x_i\le 1
\]
则每一个合取范式可以根据以下规则写成不等式$\sum_i c_i \ge 1$：
\begin{itemize}
	\item 以原变量形式出现的$x_i$以$c_i=x_i$的形式出现在不等式里
	\item 以反变量形式出现的$x_j$以$c_i=1-x_j$的形式出现在不等式里
\end{itemize}
为了变为小于号，不等式两边加上负号，有\eqref{eqn:01int_primary}
\begin{equation}
	\label{eqn:01int_primary}
	-\sum_i c_i \le -1
\end{equation}
以上过程可以在多项式时间内完成。将\eqref{eqn:01int_primary}整理成按照题目所给的形式如\eqref{eqn:01int}
\begin{equation}
	\label{eqn:01int}
	\begin{aligned}
		A_{mn}&=\left( a_{ij}=\begin{cases}
			-1 & \exists x_j \in CNF_i\\
			1 & \exists \bar{x_j} \in CNF_i\\
			0 & otherwise
		\end{cases} 
		\right) \\
		b&=\left( b_i=-1+\sum_j \mathop{\max}(0,a_{ij}) \right)
	\end{aligned}
\end{equation}
这个过程显然能在$\mathcal{O}(P(m,n))$的时间内完成。
因此该约简过程可以在多项式时间内完成。

因此，01整数规划问题是NP-Hard的。
\end{proof}

综上，01整数规划问题是NP-Complete的。

\end{spacing}

\end{document}